\section{Introduction and Positioning}
Fault assessment for a spiking controller depends on both the neural network and its feedback system. A neural-output criterion may pass before a later collision. Early stopping can change measured activity, while delayed actuation separates a logged decision from its physical consequence. Simulation engineers therefore need to check what records actually paired and observed before interpreting fault loss.

Closed-loop fault evaluation is established. DriveFI mines faults and scenarios that violate autonomous-driving safety constraints\cite{jha2019drivefi}. SpikeFI, SpikingJET and enpheeph provide configurable SNN fault campaigns\cite{spyrou2024spikefi,gogebakan2024spikingjet,colucci2022enpheeph}. CRN synchronization\cite{glasserman1992crn}, hybrid-trajectory monitoring\cite{donze2010breach} and task/resource benchmarking\cite{yik2023neurobench,dangelo2026embodied} also precede this work. Cross-layer NoC fault analysis links hardware routing defects to SNN behavior\cite{chessa2026crosslayer}; functional row deletion alone supplies no comparable hardware evidence.

We contribute a plant-neutral campaign checker, two native replay adapters, and prospective same-prior receptive-field (RF) assessments. The evaluated increment is a testable agreement about pairing, clocks, exposure and task labels, with explicit failure cases and blind spots. A delayed tank establishes limited cross-plant reuse; independent RF campaigns test the conditionality of a fly-connectome controller contrast. We do not claim a new injector, pairing estimator or general closed-loop verification method.
